\chapter{Wilson--defect parent-superconnection gate}

Reverse-audit выбрал strongest route: получить \(B-L\) vortex и family
generator \(K_n\) из одной superconnection. Кинематически это возможно.
На Nambu--family bundle можно записать
\begin{equation}
 \mathbb A=d+iq_{root}a\otimes I+I\otimes A_F+\Phi_{odd},
 \qquad
 G=U(1)_{root}\times SU(2)_F,
\end{equation}
где charge-two \(\Phi\) создаёт Majorana pairing, а adjoint action
\(A_F\) на generation triplet ограничивается на core как \(K_n/L\).
Следовательно, отдельный boundary bilinear действительно не обязателен.

\section{Почему одна запись ещё не является одним механизмом}

В минимальном local action два gauge factors остаются независимыми.
Quadratic kinetic mixing отсутствует:
\begin{equation}
 F_{root}\Tr F_F=0,
\end{equation}
поскольку \(\mathfrak{su}(2)\) traceless. На flat Wilson branch
\(F_F=0\), поэтому higher local curvature couplings также не фиксируют
holonomy. Реальный мост может возникнуть только из determinant полей,
заряженных одновременно по root и family sectors.

\section{Exact-weight obstruction}

Требуемый Wilson functional содержит
\begin{equation}
 \frac83\log(1-\cos\theta).
\end{equation}
Exact determinant означает восемь одинаковых real rotation planes веса
\(|m|=1\) и отсутствие higher weights.

Для honest \(SO(3)\) representation все spins целые. Exhaustive weight
count показывает, что единственное решение состоит из восьми vector
triplets:
\begin{equation}
 8V_1,
 \qquad
 \dim_{\mathbb R}=24.
\end{equation}
Каждый triplet содержит дополнительную zero-weight line. Их удаление
требует axis-dependent quotient и разрушает исходную covariant trace.

Альтернативный exact-log кандидат --- четыре complex fundamental doublets
\(SU(2)\), real dimension \(16\). Но half-integer representations не
спускаются к \(SO(3)\): center действует знаком \(-1\). Кроме того, при
одной и той же \(SU(2)\)-connection adjoint triplet видит удвоенный Cartan
angle. Поэтому determinant дублетов не является single-valued functional
того же \(SO(3)\) holonomy \(R_n(\theta_*)\), который создаёт core kernel.

\begin{proposition}[Representation fork]
Unbroken \(SO(3)\)-trace требует лишних zero-weight modes, а covariant
half-integer completion требует physical \(SU(2)\) lift. Однако это не
закрывает axis-broken architecture: после самостоятельного выбора оси
остаточная \(U(1)\)-ветвь может содержать восемь unit-weight rotation
planes. Поэтому representation gate является ordering condition, а не
полным no-go.
\end{proposition}

\section{Fixed-charge obstruction}

Даже four-doublet determinant закрывает только logarithmic term. Fixed-charge
Routhian требует среднего momentum norm
\begin{equation}
 \langle p^2\rangle=6.
\end{equation}
Единого quantized momentum \(|p|=\sqrt6\) не существует. Однако сохранённый
exhaustive audit восьми rotation planes даёт две положительные нечётные
ветви. При минимально возможном \(p_{\max}=3\) решение единственно:
\begin{equation}
 (p_1,\ldots,p_8)=(1,1,1,3,3,3,3,3),
 \qquad
 \sum_i p_i^2=48.
\end{equation}
При trace rank \(3\) это точно даёт
\begin{equation}
 \frac{1}{2\cdot3}\sum_i p_i^2\frac1{1-c}
 =\frac8{1-c}.
\end{equation}
Вместе с восемью unit-weight planes получается и точный logarithmic term
\(\frac83\log(1-c)\). Следовательно, оба нелинейных коэффициента действительно
выпадают из одной quantized rotor architecture без непрерывной подгонки.

Открытым остаётся не арифметика, а parent selection: почему path integral
проецируется именно на fixed-charge ensemble и почему из двух odd solutions
выбирается ветвь с минимальным максимальным momentum.

\section{Parent-action verdict}

Общая superconnection существует как совместимая кинематика. После
восстановления momentum-sector audit она также допускает exact operator-level
Wilson coefficients. Но action всё ещё содержит независимые root coupling,
pairing mass/quartic, family stiffness, Wilson tree stiffness, rotor sector
и axis coupling.

\begin{nogocriterion}[Parent-selection gate]
Точный momentum multiset нельзя считать физическим выводом, пока fixed-charge
projection, восемь unit-weight planes и criterion выбора minimal
\(p_{\max}\) не следуют из заранее объявленной boundary path integral.
Использование этих данных только потому, что они дают коэффициенты, было бы
дискретной подгонкой.
\end{nogocriterion}

Таким образом, мой первоначальный verdict о полном закрытии был чрезмерным.
Wilson--defect bridge переоткрыт условно: два нелинейных коэффициента уже
получены точно на operator level. Следующий decisive расчёт --- canonical
boundary path integral с projector на fixed momentum, включая measure,
zero modes, BV/BRST content и coupling к family axis.

\section{Family-Casimir selector вместо
\texorpdfstring{minimal-\(p_{\max}\)}{minimal-pmax}}

Generation triplet даёт каноническое разложение
\begin{equation}
 \operatorname{End}_0(\mathbb R^3)
 =\Lambda^2\mathbb R^3\oplus\operatorname{Sym}^2_0\mathbb R^3
 \simeq V_1\oplus V_2,
 \qquad 8=3+5.
\end{equation}
На блоках \(V_1,V_2\) family-Casimir равен \(2,6\). Поэтому оператор
\begin{equation}
 P=\frac12C_2
\end{equation}
имеет spectrum \(1^3 3^5\) и
\begin{equation}
 \Tr P^2=3\cdot1^2+5\cdot3^2=48.
\end{equation}
Альтернативная ветвь с momentum \(5\) потребовала бы Casimir eigenvalue
\(10\), которого в \(V_0\oplus V_1\oplus V_2\) нет. Это доказывает наличие
family-оператора с нужным spectrum, но ещё не выбор rotor state:
\(C_2/2\) и \(p=-i\partial_\phi\) действуют на разных сомножителях.

\section{Exact reduced path integral}

Однослойный rotor no-go не запрещает две градации одной boundary
architecture. Compact grade имеет rank \(8\) и fixed sector \(P=C_2/2\), а
unit-root gaussian grade равен
\begin{equation}
 \mathbb R^2_{\rm root}\otimes\operatorname{End}_0(\mathbb R^3),
 \qquad \dim_{\mathbb R}=16.
\end{equation}
При общей inertia \(I=1-c\)
\begin{equation}
 Z_P=\exp\left(-\frac{24}{I}\right),
 \qquad
 Z_G\propto I^{-8}.
\end{equation}
Поэтому после объявления free energy на одну family component
\begin{equation}
 -\frac13\log(Z_PZ_G)
 =\frac8{1-c}+\frac83\log(1-c)+\text{const}.
\end{equation}
Оба nonlinear terms воспроизводятся обычными integer field counts;
fractional determinant не требуется. Однако fixed-charge projector и общая
family normalization должны быть выведены отдельно.

\section{Mapping-cone common operator}

Общая inertia сама выводится из Wilson incidence. На unit-root plane
\begin{equation}
 U(\theta)=
 \begin{pmatrix}c&-s\\s&c\end{pmatrix},
 \qquad Q_U=1-U(\theta),
\end{equation}
и при \(c^2+s^2=1\)
\begin{equation}
 \frac12Q_U^TQ_U=(1-c)I_2.
\end{equation}
Один mapping-cone Laplacian одновременно задаёт rotor inertia и gaussian
mass. Он действует на root factor и коммутирует с family-Casimir, поэтому
блоки \(V_1\) и \(V_2\) не смешиваются на quadratic level.

Canonical Wilson plaquette \(1-c\) даёт требуемую unit stiffness \(-c\) с
точностью до константы. В результате единая quadratic architecture выводит
\begin{equation}
 V_{\rm cone}(c)
 =\frac8{1-c}+\frac83\log(1-c)-c.
\end{equation}
От полного функционала остаётся ровно
\begin{equation}
 V_{\rm gap}(c)-V_{\rm cone}(c)=\frac{c}{45}.
\end{equation}

\begin{proposition}[Mapping-cone pass]
Family selector, nonlinear pair, их общая inertia и canonical \(-c\)
stiffness теперь происходят из одной локальной quadratic architecture.
\end{proposition}

Число \(1/45\) уже известно как constant response периодического phase-zero
channel, но его primitive \(+c/45\) ещё не встроен в тот же local BV
complex.

\begin{nogocriterion}[Periodic-singlet embedding gate]
Полное Wilson closure допускается только если phase-zero singlet с response
\(1/45\) выводится из той же superconnection с фиксированными statistics,
zero-mode prescription и determinant sign. Ручное добавление \(+c/45\)
запрещено.
\end{nogocriterion}

Следующий расчёт теперь предельно узок: встроить один periodic singlet и
проверить полный БВ-якобиан.

\section{Three-node graded completion}

Независимый parity/charge audit показывает, что two-node bundle не может
одновременно содержать unit-charged zero-form и unit-charged one-form полной
нечётной степени. Минимальный carrier имеет три summands:
\begin{equation}
 E=E_0^+\oplus E_1^-\oplus E_2^+,
 \qquad (q_0,q_1,q_2)=(0,1,1).
\end{equation}
В нём zero-form \(E_0^+\to E_1^-\) и one-form \(E_0^+\to E_2^+\)
имеют unit root charge и видят один оператор
\begin{equation}
 K_\theta=\frac{2-U_\theta-U_\theta^\dagger}{2}=1-c.
\end{equation}
Quadratic cross term исчезает по form degree и ортогональности targets.
После compact reduction one-form grade даёт восемь angular rotors, а
zero-form grade --- шестнадцать gaussian modes.

Вместе с canonical plaquette и periodic primitive reduced action равен
\begin{equation}
 \Gamma_{\rm red}(c)
 =\frac8{1-c}+\frac83\log(1-c)+(1-c)+\frac{c}{45},
\end{equation}
то есть совпадает с \(V_{\rm gap}\) с точностью до константы.

\begin{proposition}[Three-node quadratic embedding pass]
Трёхузловая graded superconnection является минимальным carrier обеих
градаций, общего Wilson operator и полной reduced shape без нового
continuous polarization.
\end{proposition}

Это quadratic configuration-metric pass, а не final physical closure.
Supercurvature interactions начинаются выше quadratic order; periodic
primitive, compact radial constraint, gauge quotient и interaction Jacobian
должны быть вычислены из одного action.

\begin{nogocriterion}[Curvature/BV gate]
Three-node branch проходит только если одно заранее фиксированное curvature
или spectral action после gauge fixing сохраняет common \(K_\theta\), Casimir
projector, sixteen-mode determinant и linear primitive без counterterm.
\end{nogocriterion}

Машинный parity/charge audit сохранён в
\path{s2t_v4_three_node_superconnection_closure_gate.py} и
\path{s2t_v4_three_node_superconnection_closure_gate_results.json}.

\section{Periodic-singlet determinant: all-orders BV no-go}

Положим \(t=1-c\). Для periodic nonzero tower normalized determinant имеет
разложение
\begin{equation}
 \sum_{n\ne0}\log\left(1+\frac{t}{(\pi n)^4}\right)
 =\frac{t}{45}-\frac{t^2}{9450}+O(t^3),
\end{equation}
поскольку
\begin{equation}
 \frac{2\zeta(4)}{\pi^4}=\frac1{45},
 \qquad
 \frac{2\zeta(8)}{\pi^8}=\frac1{4725}.
\end{equation}
Fermionic/ghost sign даёт требуемый первый член \(-t/45\), эквивалентный
\(+c/45\) с точностью до константы, но одновременно оставляет лишний
\(+t^2/9450\). Поэтому susceptibility coefficient нельзя выдавать за полный
determinant.

Для конечного BV menu с couplings \(\lambda_j\) и signed weights \(a_j\)
точная линейность потребовала бы
\begin{equation}
 \sum_j a_j\lambda_j^k=0,
 \qquad k\ge2.
\end{equation}
При \(m\) различных ненулевых couplings условия \(k=2,\ldots,m+1\)
образуют невырожденную weighted Vandermonde system. Значит, все net weights
равны нулю и linear moment также исчезает.

\begin{theorem}[Finite Gaussian periodic-singlet no-go]
Никакой конечный стандартный boson/fermion/ghost Gaussian BV complex не даёт
ненулевой exact linear primitive \(+c/45\) при сокращении всех higher powers.
\end{theorem}

Остаётся иной класс:
\begin{equation}
 \Gamma_\zeta(t)
 =-\frac{t}{\pi^4}\Tr{}'\Delta_{\rm per}^{-2}
 =-\frac{t}{45}.
\end{equation}
Это exact nonlocal spectral stiffness, а не determinant конечного local
Gaussian complex.

\begin{nogocriterion}[Zeta-stiffness reopening]
Ветвь переоткрывается только если этот linear zeta functional с unit
coefficient и отрицательным знаком выводится из parent spectral action,
topological expectation value или infinite complex с доказанным all-orders
cancellation. Ручной counterterm запрещён.
\end{nogocriterion}

\section{Periodic exact--ghost BV ledger}

В полном gauge-fixed учёте нельзя оставлять один ghost без его продольного
bosonic partner. Для одной real root plane имеются две вещественные gauge
generators. Exact one-form sector даёт signed Gaussian weight
\begin{equation}
 +\frac12\times2=+1,
\end{equation}
а две соответствующие FP ghost pairs дают
\begin{equation}
 -1\times2=-2.
\end{equation}
Их полный net weight равен
\begin{equation}
 +1-2=-1.
\end{equation}
Поэтому periodic trace автоматически создаёт
\begin{equation}
 \Gamma_{\rm ex+gh}(c)
 =-(1-c)\frac{2\zeta(4)}{\pi^4}
 =-\frac{1-c}{45}
 =\frac{c}{45}+\text{const}.
\end{equation}
Это ровно недостающий term, уже с полным exact/ghost bookkeeping, а не с
изолированным ghost.

\begin{proposition}[Periodic BV ledger pass]
Root-plane exact one-form и его FP ghosts имеют канонический net weight
\(-1\), который вместе с periodic \(n^{-4}\) trace даёт exact correction
\(-(1-c)/45\).
\end{proposition}

Обычный четырёхмерный Yang--Mills complex, однако, содержит также coexact
one-form modes с положительным bosonic weight. Они добавили бы новые spectral
terms. Поэтому pass становится физическим только для topological/flat
one-form grade либо при обязательном BV partner, сокращающем весь coexact
root sector.

\begin{nogocriterion}[Coexact-cancellation gate]
Ветвь проходит только если local three-node action выводит отсутствие или
точное BV-сокращение coexact root modes, сохраняя net exact/ghost weight
\(-1\) и compact harmonic rotor.
\end{nogocriterion}

\section{Coexact-cancellation trilemma}

Полный Hodge/BV audit оставляет три стандартные возможности.
\begin{enumerate}
 \item Ordinary Maxwell-type kinetic term сохраняет положительный coexact
 root tower. Он не сокращается уже использованными FP ghosts.
 \item Pure flat/BF constraint удаляет coexact modes, но одновременно
 превращает exact/ghost sector в topological torsion и harmonic measure;
 требуемый \(-\Tr'\Delta^{-2}\) residue не сохраняется.
 \item Dynamical или localized BF partner меняет eigenvalues либо boundary
 conditions и потому не является selective isospectral cancellation.
\end{enumerate}

\begin{theorem}[Local one-form BV trilemma]
В стандартном конечном local one-form/BF/BV class нельзя одновременно
сохранить periodic exact--ghost stiffness \(-(1-c)/45\) и удалить все
coexact root modes. Ordinary kinetic branch оставляет лишний tower, а flat
topological branch удаляет также нужный spectral residue.
\end{theorem}

Это закрывает не всю Wilson reconstruction, а её local Gaussian one-form
реализацию. Остаётся nonlocal spectral class
\begin{equation}
 \Gamma_{\rm stiff}(c)
 =-\frac{1-c}{\pi^4}\Tr{}'\Delta_{\rm per}^{-2}
 =-\frac{1-c}{45},
\end{equation}
прочитанный как fundamental Kubo/configuration metric, а не как determinant
локального finite BV complex.

\begin{nogocriterion}[Fundamental spectral-metric gate]
Ветвь переоткрывается только если unit coefficient, negative supertrace sign
и kernel \(\Delta^{-2}\) выводятся из parent information/spectral geometry до
сравнения с Wilson angle. Объявление этого functional вручную запрещено.
\end{nogocriterion}

Машинный аудит сохранён в
\path{s2t_v4_periodic_singlet_bv_embedding_gate.py} и
\path{s2t_v4_periodic_singlet_bv_embedding_gate_results.json}.

Машинный exhaustive audit сохранён в
\path{s2t_v4_wilson_defect_parent_superconnection_gate.py} и
\path{s2t_v4_wilson_defect_parent_superconnection_gate_results.json}.

\section{Fixed-charge Gauss-law и общая family-мера}

Следующий аудит разделяет две ранее склеенные задачи: происхождение деления
на три и динамический выбор momentum vector
\begin{equation}
 p_*=(1,1,1,3,3,3,3,3).
\end{equation}

\subsection{Нормировку на три можно сделать согласованной}

Проблема деления \(-\log Z\) на три исчезает, если сначала строится полный
family-traced functional, в котором tree и periodic sectors также имеют
triplet multiplicity. Тогда
\begin{equation}
 \Gamma_{\rm full}(c)
 =\frac{24}{1-c}+8\log(1-c)-3c+\frac{c}{15}
 =3V_{\rm gap}(c).
\end{equation}
Деление на три после этой формулы является только переходом к величине на
одну family component и не меняет stationary equation. Это не полный
динамический вывод меры, но устраняет прежнюю внутреннюю несогласованность,
где determinant делился на три, а tree term --- нет.

\subsection{Один закон Гаусса не выбирает восемь импульсов}

Пусть восемь compact momenta образуют vector \(p\in\mathbb Z^8\). Один
temporal gauge field вводит одно scalar Gauss constraint
\begin{equation}
 A p=k,
 \qquad \rank A=1.
\end{equation}
Его solution lattice имеет dimension $7$, поэтому он не может единственным
образом зафиксировать $p=p_*$. Даже два block constraints, различающие
$V_1$ и $V_2$, имеют rank $2$ и оставляют шестимерное семейство.
Уникальное линейное закрепление всего vector требует rank $8$:
\begin{equation}
 p-C_2/2=0
\end{equation}
как восьмикомпонентного operator constraint. Такой multiplier не следует из
одного минимального $U(1)$; он является новой структурой, пока не выведен
из существующей superconnection.

Есть и второе различие. Projector в Hilbert space одного rotor,
тензорно умноженного на family carrier, даёт сумму
\begin{equation}
 Z_{\rm dir}=3e^{-1/[2(1-c)]}+5e^{-9/[2(1-c)]},
\end{equation}
а использованная ранее формула
\begin{equation}
 Z_{\rm prod}=e^{-24/(1-c)}
\end{equation}
является произведением восьми независимо constrained rotors. Поэтому один
operator projector не восстанавливает прежний path integral автоматически.

\begin{proposition}[Minimal Gauss-law menu no-go]
Один diagonal $U(1)$ Gauss law и два irrep-block Gauss laws не выводят
единственный sector $1^3 3^5$. Family normalization допускает согласованное
triplet-trace completion, но fixed-charge product требует дополнительного
rank-eight covariant constraint.
\end{proposition}

Следующий честный gate имеет два исхода. Либо rank-eight constraint выводится
из уже существующей graded superconnection без нового свободного multiplier,
либо восьмироторная реконструкция снимается и заменяется другим механизмом
inverse term.

\section{Coherent-source Gaussian: альтернативная роторная конструкция}

Более широкий пересмотр показывает, что inverse term необязательно должен
происходить из fixed momentum. Пусть уже имеющийся Gaussian carrier равен
\begin{equation}
 \mathcal X=\mathbb R^2_{\rm root}\otimes(V_1\oplus V_2),
 \qquad \dim_{\mathbb R}\mathcal X=16.
\end{equation}
Для выбранной family axis $n$ обозначим через $\kappa_j(n)$ нормированный
reproducing-kernel vector spin-$j$ channel. При нормировке invariant measure
его квадрат нормы равен размерности representation:
\begin{equation}
 \|\kappa_1(n)\|^2=3,
 \qquad
 \|\kappa_2(n)\|^2=5.
\end{equation}
Это не выбор восьми компонент вручную: при вращении axis vector
$\kappa_j(n)$ преобразуется ковариантно, а его норма остаётся постоянной.

Положим
\begin{equation}
 P=\frac12C_2,
 \qquad
 J(n)=e_{\rm root}\otimes P\kappa(n).
\end{equation}
Тогда representation theory автоматически даёт
\begin{equation}
 \|J(n)\|^2
 =3\cdot1^2+5\cdot3^2
 =48,
\end{equation}
независимо от axis. Рассмотрим один Gaussian integral с общей mapping-cone
inertia $I=1-c$:
\begin{equation}
 Z_J(I)=\int_{\mathbb R^{16}}dX\,
 \exp\left[-\frac I2\|X\|^2+i\langle J(n),X\rangle\right].
\end{equation}
Точное completing the square даёт
\begin{equation}
 Z_J(I)\propto I^{-8}\exp\left[-\frac{24}{I}\right].
\end{equation}
Следовательно,
\begin{equation}
 -\frac13\log Z_J
 =\frac8{1-c}+\frac83\log(1-c)+\text{const}.
\end{equation}

Таким образом, один и тот же sixteen-mode Gaussian sector порождает оба
nonlinear terms. Восемь compact rotors, momentum vector $1^3 3^5$ и
rank-eight Gauss projector больше не нужны. Casimir теперь задаёт norm
ковариантного source, а не навязанные quantum numbers.

\begin{proposition}[Coherent-source algebraic pass]
Axis-coherent source $J(n)=e_{\rm root}\otimes(C_2/2)\kappa(n)$ в
sixteen-dimensional Gaussian carrier точно воспроизводит nonlinear pair
$8/(1-c)+(8/3)\log(1-c)$ без fixed-charge sector.
\end{proposition}

Это существенное переоткрытие, но не физическое замыкание. Для положительного
inverse term требуется imaginary Euclidean coupling
$i\langle J,X\rangle$. Он допустим как кандидат на oriented defect или
boundary phase, однако его gauge-invariant происхождение, contour,
reflection/reality prescription и measure axis должны следовать из parent
superconnection. Вставка source вручную снова была бы подгонкой.

\begin{nogocriterion}[Coherent-source parent gate]
Ветвь проходит только если одна заранее фиксированная defect coupling выводит
imaginary source, его unit coefficient, root orientation и invariant
coherent-state measure. Сравнение с $V_{\rm gap}$ проводится только после
фиксации этих данных.
\end{nogocriterion}